\documentclass{article}
\usepackage{iclr2026_conference,times}
\input{math_commands.tex}

\usepackage{hyperref}
\usepackage{url}
\usepackage{graphicx}
\usepackage{booktabs}
\usepackage{amsmath}
\usepackage{amssymb}
\usepackage{xcolor}

\newcommand{\figOverviewPath}{../figures/adrd_overview.pdf}
\newcommand{\figBudgetFrontierPath}{../figures/budget_frontier.pdf}
\newcommand{\figSameBudgetPath}{../figures/same_budget_comparison.pdf}
\newcommand{\figLongContextPath}{../figures/long_context.pdf}
\newcommand{\tableMainResultsPath}{../figures/main_results_table.tex}
\newcommand{\tableModeUsagePath}{../figures/mode_usage_table.tex}


\newcommand{\method}{ADRD}
\newcommand{\factorized}{\textsc{Factorized}}
\newcommand{\coupled}{\textsc{Coupled-Local}}
\newcommand{\micro}{\textsc{Micro-Sequential}}

\title{Budgeted Selective Dependency Routing for Discrete Diffusion Language Models\\A Synthetic Feasibility Study}

\author{}

\begin{document}
\maketitle

\begin{abstract}
Discrete diffusion language models offer attractive parallel decoding, but their quality can degrade when simultaneously predicted tokens have strong mutual dependencies. Recent work typically responds by choosing one dependency regime globally: stay factorized, introduce a stronger coupled output family, or move toward more sequential decoding. We study a narrower question: can richer dependency handling be allocated only to the minority of regions that appear risky? We present Adaptive Dependency Routing Diffusion (\method), a budgeted controller that routes unresolved regions to one of three decoding modes: factorized prediction, short-span coupled decoding, or tiny micro-sequential updates. The current implementation is intentionally modest: a synthetic discrete-diffusion scaffold with heuristic routing rather than a trained language model. Even in that setting, selective routing improves token accuracy from $0.6222$ to $0.6825$ and dependency accuracy from $0.3446$ to $0.6245$ while spending extra computation on only $10.89\%$ of tokens. On a longer-context stress test, \method\ improves dependency accuracy by $12.95$ points. At the same time, a simpler coupled-only surrogate slightly outperforms full \method\ on token accuracy under the relaxed budget, so our evidence supports feasibility rather than a definitive adaptive-routing win. The resulting paper is a controlled pilot study: selective dependency allocation looks promising, but stronger validation on real diffusion language models remains necessary.
\end{abstract}

\section{Introduction}

Diffusion language models are attractive because they can update many tokens in parallel, but that same parallelism creates a structural problem: simultaneously decoded tokens are often treated as conditionally independent even when the sequence requires coordinated decisions. The resulting factorization gap is now a central bottleneck for discrete diffusion decoding. Recent work addresses the problem in several different ways, including stronger coupled output layers \citep{li2026codd}, more efficient block or semi-autoregressive decoding \citep{liu2026sdlm,wu2026fastdllmv2}, inference-time remasking or self-correction \citep{wang2025remasking,zhang2026selfcorrection}, and broader autoregressive-diffusion hybrids \citep{fathi2025unifying,sahoo2026esolm}. What remains less explored is whether the expensive dependency model should be chosen \emph{locally} rather than globally.

This paper studies that question in deliberately constrained form. Instead of claiming a new state-of-the-art diffusion language model, we isolate the controller hypothesis on a small synthetic scaffold that exposes copy-like and progression-like dependency groups. Our proposal, Adaptive Dependency Routing Diffusion (\method), routes each unresolved region at each denoising step to one of three modes: \factorized\ decoding for cheap independent updates, \coupled\ decoding for short strongly linked spans, and \micro\ decoding for tiny high-risk knots. A budget controller constrains how many tokens may receive non-factorized treatment. The design target is not ``joint modeling everywhere'' but ``pay for richer dependency handling only where the factorization gap appears largest.''

The resulting evidence is useful but mixed. Under this local scaffold, selective routing clearly improves over a fully factorized baseline. At a relaxed span-4 budget, \method\ increases token accuracy from $0.6222$ to $0.6825$ and dependency accuracy from $0.3446$ to $0.6245$ while using expensive routing for only $10.89\%$ of tokens. It also improves a longer-context stress setting. However, the strongest same-budget surrogate is not the full adaptive policy: a \coupled-only baseline slightly exceeds \method\ on token accuracy, and a random router is competitive on some dependency metrics. This means the present paper supports a feasibility claim, not a final claim that adaptive three-way routing is already the best policy.

Our contribution is therefore a controlled pilot study with three parts:
\begin{enumerate}
\item We formulate the \method\ routing problem as selective dependency allocation under explicit token and group budgets.
\item We implement a minimal synthetic evaluation scaffold that measures token accuracy, dependency recovery, mode usage, and matched-budget tradeoffs.
\item We show that selective non-factorized routing is viable, while also identifying the main failure mode: on the current scaffold, most gains can already be explained by a simpler \coupled-only policy.
\end{enumerate}

This partly negative outcome is a feature rather than an embarrassment. A clean pilot paper should tell the reader not only that selective routing can help, but also which part of the design currently carries the improvement and which part still lacks convincing evidence.

The rest of the paper is organized as follows. Section~\ref{sec:related} positions \method\ relative to recent diffusion decoding work. Section~\ref{sec:method} defines the routing mechanism and budget controller. Section~\ref{sec:setup} describes the synthetic task and matched-budget protocol. Section~\ref{sec:results} presents the frontier, ablations, and long-context stress test. Section~\ref{sec:limitations} explains why the current evidence should be read as synthetic feasibility rather than a finished dLLM result.

\section{Related Work}
\label{sec:related}

\paragraph{Discrete diffusion language modeling.}
Discrete diffusion models extend denoising diffusion ideas to categorical state spaces \citep{austin2021d3pm}. For text, diffusion-style language modeling has been explored in continuous embedding spaces \citep{strudel2022selfconditioned} and, more recently, in masked or discrete token spaces that emphasize parallel generation. These models motivate the core appeal of diffusion for language: generation can in principle be faster than autoregressive decoding when many tokens are updated simultaneously.

\paragraph{The factorization barrier.}
Parallel decoding becomes difficult when simultaneously updated tokens are strongly dependent. Recent work argues that this limitation is structural rather than incidental: the output family is often fully factorized even when the task requires joint reasoning. CoDD is the clearest recent statement of this argument, replacing factorized outputs with a tractable coupled inference layer \citep{li2026codd}. Our work differs in emphasis. We do not propose a better global output family. Instead, we ask whether a cheaper model can remain factorized on easy regions while escalating only a small subset of spans to richer handlers.

\paragraph{Efficient diffusion decoding.}
Several lines of work try to make diffusion decoding practical without fully abandoning parallelism. Sequential Diffusion Language Models move toward dynamic subsequence decoding \citep{liu2026sdlm}. Fast-dLLM v2 uses block diffusion and caching to improve the speed-quality frontier of diffusion LLMs \citep{wu2026fastdllmv2}. These methods show that partial sequential structure and systems-aware decoding can be valuable. \method\ is complementary: it focuses on \emph{where} extra dependency handling should be applied under a fixed budget.

\paragraph{Inference-time repair and hybridization.}
Remasking and self-correction approaches spend additional inference compute after an initial diffusion update \citep{wang2025remasking,zhang2026selfcorrection}. Hybrid AR/diffusion models further blur the line between full parallelism and sequential generation \citep{fathi2025unifying,sahoo2026esolm}. Those works motivate our controller perspective. If extra compute already helps, then the next question is whether a routing policy can allocate that compute more selectively to the regions that need it most.

\paragraph{Position of this paper.}
The present paper should be read as a feasibility probe that sits between global coupling and sampler heuristics. It does not claim to replace CoDD-style joint modeling, nor to compete directly with production dLLMs. Its narrower contribution is to show that selective dependency allocation is measurable, budgetable, and diagnostically interpretable in a controlled setting.

\section{Adaptive Dependency Routing Diffusion}
\label{sec:method}

We consider a masked sequence at diffusion step $t$ with a set of unresolved regions $\mathcal{R}_t$. Each region is a short subset of token positions that appears mutually dependent according to a dependency estimator. In the local implementation, these regions are synthetic copy or progression groups exposed by the benchmark generator, but the routing abstraction itself is broader.

\begin{figure}[t]
\centering
\includegraphics[width=\linewidth]{\figOverviewPath}
\caption{Overview of \method. A dependency estimator scores unresolved regions, a budget controller restricts the use of expensive modes, and the router assigns each region to the cheapest mode that appears adequate.}
\label{fig:overview}
\end{figure}

\paragraph{Three decoding modes.}
Each unresolved region $r \in \mathcal{R}_t$ is assigned to one of three modes:
\begin{itemize}
\item \factorized: predict tokens independently and update the region in one cheap parallel step.
\item \coupled: apply a short-span joint handler to a small local group, intended for copy-like or strongly coupled spans.
\item \micro: decode a tiny subset sequentially, then return to parallel denoising.
\end{itemize}
The key idea is not that one of these modes is universally best, but that different regions may need different levels of dependency modeling.

\paragraph{Dependency estimator.}
For each region $r$, the estimator returns a risk score $s_t(r)$ and a confidence proxy $c_t(r)$. The current code uses a heuristic version of this estimator in which larger unresolved spans and progression-like patterns receive higher risk scores. This choice keeps the experiment transparent: routing gains cannot be attributed to a hidden high-capacity learned model.

\paragraph{Budgeted router.}
The router selects a mode $m_t(r) \in \{F,C,S\}$ for factorized, coupled-local, or micro-sequential decoding. Let $|r|$ denote the number of unresolved tokens in region $r$. At each diffusion step we constrain the total amount of expensive routing by
\begin{align}
\sum_{r \in \mathcal{R}_t} \mathbf{1}[m_t(r) \neq \factorized] |r| &\leq B_{\mathrm{tok}}, \\
\sum_{r \in \mathcal{R}_t} \mathbf{1}[m_t(r) = \coupled] &\leq B_{\mathrm{cpl}}, \\
\sum_{r \in \mathcal{R}_t} \mathbf{1}[m_t(r) = \micro] &\leq B_{\mathrm{mic}}.
\end{align}
The implementation additionally caps the maximum size of a \coupled\ span. This prevents the method from quietly degenerating into broad joint decoding.

\paragraph{Routing policy.}
The current policy is heuristic. Progression-like groups and very small high-risk groups are biased toward \micro; copy-like groups and larger short spans are biased toward \coupled. If neither expensive option fits the remaining budget, the region falls back to \factorized. This design is intentionally simple because the paper's goal is to test whether selective allocation is worth pursuing at all.

\paragraph{Local decoding operators.}
The three local operators are also synthetic. \factorized\ predictions are noisy and independent. \coupled\ decoding is more reliable on short groups, while \micro\ decoding first predicts an anchor token with high confidence and then conditionally resolves the rest of the region. These operators are not intended as faithful substitutes for a production dLLM; they serve as controlled surrogates for cheap, local-joint, and local-sequential behavior.

\paragraph{Why this formulation matters.}
The formulation forces the paper to answer a specific question: if richer dependency models are expensive, can a controller spend that budget on the right minority of regions? That question remains meaningful even though the present implementation is synthetic and heuristic.

\section{Synthetic Evaluation Protocol}
\label{sec:setup}

\paragraph{Task generator.}
The benchmark is intentionally small and fully controlled. Each example is a discrete sequence generated by the local scaffold in the repository's synthetic data module. Tokens are sampled uniformly, then one or two dependency groups may be injected. A \emph{copy} group forces a span to share the same token. A \emph{progression} group enforces an arithmetic-style progression modulo the vocabulary size. These two patterns are simple, but they make cross-token dependency structure explicit and measurable.

\paragraph{Default setting.}
The main experiments use 192 sequences of length 24, vocabulary size 10, and four diffusion steps. Dependency groups appear with probability 0.8 and have span lengths between 2 and 5 tokens. The long-context stress test increases the sequence length to 40, the dependency probability to 0.9, and the span length range to 3--7.

\paragraph{Metrics.}
We report overall token accuracy, dependency-token accuracy, dependency-group accuracy, non-dependency accuracy, and the average fraction of tokens that receive non-factorized treatment. Dependency-group accuracy is especially important because a region is counted as correct only if every token in the group is recovered correctly. This metric penalizes partial recovery of coupled spans.

\paragraph{Systems compared.}
Our main comparisons are:
\begin{itemize}
\item a fully factorized baseline;
\item strict-budget \method\ runs at 5\% and 10\% non-factorized token budgets;
\item a relaxed span-4 \method\ run with a 20\% cap;
\item matched-budget fixed surrogates that allow only \coupled\ or only \micro;
\item and a random matched-budget router.
\end{itemize}

\paragraph{Budget discipline.}
Every comparison in Section~\ref{sec:results} is budget matched. The method is not allowed to claim gains by simply spending unbounded extra compute. Under the strict 10\% setting, the actual non-factorized token fraction is only 2.93\%; under the relaxed span-4 setting it is 10.89\%.
For the relaxed-budget ablations, the \coupled-only, \micro-only, random, and full \method\ systems all operate under the same top-level budget cap, which makes the comparison about allocation policy rather than raw compute.

\begin{table}[t]
\centering
\caption{Main synthetic runs. The two-mode 10\% system already improves over fully factorized decoding, and the relaxed span-4 \method\ run improves further while still using non-factorized handling on a minority of tokens.}
\label{tab:main}
\input{\tableMainResultsPath}
\end{table}

\paragraph{Scope.}
This protocol is intentionally narrower than a real diffusion language modeling benchmark. There is no pretrained backbone, no human evaluation, and no external latency measurement. The benefit is that the routing signal, the compute budget, and the dependency-heavy regions are all directly inspectable.

\section{Results}
\label{sec:results}

\subsection{Budgeted routing beats the fully factorized baseline}

Table~\ref{tab:main} and Figure~\ref{fig:frontier} show the first clear result: once the controller is allowed to spend a small amount of extra budget, the synthetic factorization gap shrinks substantially. Moving from the factorized baseline to the strict 10\% configuration improves token accuracy from $0.6222$ to $0.6413$ and dependency accuracy from $0.3446$ to $0.4363$. Relaxing the budget further to the span-4 configuration improves token accuracy to $0.6825$ and dependency accuracy to $0.6245$.

The 5\% budget setting is also informative. In that regime, the controller collapses back to fully factorized behavior and performance is unchanged. This indicates that the gains are not a free side effect of the router logic; they appear only when the controller is permitted to spend enough budget to activate richer dependency handling.

\begin{figure}[t]
\centering
\includegraphics[width=\linewidth]{\figBudgetFrontierPath}
\caption{Budget frontier on the synthetic task. A very small budget collapses to factorized behavior, while a relaxed but still minority budget produces large dependency gains.}
\label{fig:frontier}
\end{figure}

\subsection{The current adaptive policy is viable, but not yet dominant}

The stronger question is whether the adaptive three-mode policy is itself superior to simpler matched-budget alternatives. Figure~\ref{fig:matched} shows that the answer is not yet yes. Under the relaxed span-4 budget, the \coupled-only surrogate attains the best token accuracy ($0.6871$), slightly above \method\ ($0.6825$). The random router reaches the highest dependency accuracy on this specific synthetic slice, while \method\ stays near the top without clearly dominating any single metric. Full \method\ remains competitive and uses both expensive modes, but the evidence does not justify a claim that adaptive routing is already the best same-budget strategy.

\begin{figure}[t]
\centering
\includegraphics[width=\linewidth]{\figSameBudgetPath}
\caption{Matched-budget comparison under the relaxed span-4 setting. \method\ is competitive and mixes both expensive modes, but a simpler \coupled-only surrogate is slightly stronger on token accuracy.}
\label{fig:matched}
\end{figure}

That negative result is scientifically useful. It narrows where future work should go. The router is not yet the star of the system; the main gains seem to come from allowing \coupled\ decoding at all. A stronger paper would need either a learned router that beats fixed policies under the same budget, or a richer task where \micro\ updates contribute more clearly and more often.

\begin{table}[t]
\centering
\caption{Mode usage for the relaxed span-4 \method\ run. The controller uses both expensive modes, which means the negative result is not caused by complete collapse to a single expensive option.}
\label{tab:modes}
\input{\tableModeUsagePath}
\end{table}

Table~\ref{tab:modes} confirms that the relaxed-budget run does not completely collapse to one mode. The controller routes 407 tokens to \coupled\ and 95 tokens to \micro. That is enough to support a narrow interpretability claim: the controller can express heterogeneous mode usage. It is not enough to support the stronger performance claim reviewers would want from a main empirical paper.

\subsection{Longer contexts reinforce the same conclusion}

The length-40 stress test in Figure~\ref{fig:long} strengthens the synthetic feasibility story. \method\ improves token accuracy from $0.6268$ to $0.6526$ and dependency accuracy from $0.3720$ to $0.5015$. The absolute dependency gain of 12.95 points is one of the clearest positive results in the study.

\begin{figure}[t]
\centering
\includegraphics[width=0.82\linewidth]{\figLongContextPath}
\caption{Long-context stress test with length-40 sequences. \method\ retains a directional advantage over the factorized baseline.}
\label{fig:long}
\end{figure}

At the same time, the long-context run reinforces the earlier caution: almost all of its expensive routing goes to \coupled\ rather than \micro. The takeaway is consistent across experiments. Selective dependency handling helps, but the present scaffold still behaves as though local coupling is the main useful expensive operation.

\subsection{Interpretation}

Taken together, the experiments support three statements. First, the factorization gap is heterogeneous enough that non-factorized computation can be concentrated on a minority of tokens. Second, there is a real budget frontier rather than a binary choice between full factorization and full coupling. Third, the current heuristic controller is not yet the final answer: simpler same-budget policies can match or exceed it depending on the metric. This is why we frame the paper as a feasibility study rather than a finished adaptive-routing result.

\section{Limitations And Next Steps}
\label{sec:limitations}

This paper has several important limitations.

\paragraph{Synthetic task only.}
All evidence comes from a local synthetic benchmark. The task is useful for isolating dependency structure, but it is not a substitute for evaluation on real text, real denoisers, or real latency-sensitive decoding.

\paragraph{Heuristic routing.}
The current estimator and router are not learned. They are hand-designed rules that expose the controller hypothesis as cleanly as possible. This makes the experiment interpretable, but it also means the present paper cannot claim that a trainable routing network has been validated.

\paragraph{No direct CoDD implementation.}
Although the paper is motivated by the factorization barrier literature, this repo does not contain a real CoDD-style baseline. The ``coupled-only'' surrogate is only a synthetic local analogue. A stronger follow-up paper would compare against an actual globally coupled dLLM implementation under matched compute.

\paragraph{No external repair baseline.}
The workspace does not contain a full remasking or self-correction implementation. Those baselines are cited for context but not reproduced here.

\paragraph{What would upgrade the paper.}
The clearest next step is to port the routing abstraction to a public dLLM backbone and ask the same question under fixed denoiser calls or matched wall-clock time. The router should be learned, not heuristic, and the evaluation should include direct comparisons against global coupling, remasking, and hybrid AR/diffusion baselines. Until then, the current result should be read as a scoped but honest pilot.

\section{Conclusion}

We studied whether the factorization gap in discrete diffusion decoding can be addressed selectively rather than globally. On a controlled synthetic scaffold, Adaptive Dependency Routing Diffusion improves over fully factorized decoding while using non-factorized handling on only a minority of tokens. The gains are especially clear on dependency-sensitive and longer-context settings. However, the best current same-budget surrogate is not the full adaptive policy but a simpler coupled-only alternative. That negative result matters: it shows that selective dependency allocation is worth studying, but also that the controller itself still needs stronger justification.

The most defensible conclusion is therefore modest. Selective dependency handling appears feasible and measurable. A budgeted routing view of diffusion decoding is plausible. But the present evidence is still a pilot study, and the next serious test must happen on real diffusion language models.


\bibliographystyle{iclr2026_conference}
\bibliography{references}

\end{document}
